
\section*{5. Homogeneous spaces of infinitesimal formal groups over a field}
\phantomsection
\label{sga3:VIIB:5}
\addcontentsline{toc}{section}{5. Homogeneous spaces of infinitesimal formal groups over a field}

\subsection*{5.0}
\phantomsection
\label{sga3:VIIB:5.0}

For simplicity, suppose that \(k\) is a field.\footnote{%
\label{sga3:VIIB:nde:142}The editors added subsection 5.0 in order
to express \SGARefLink{sga3:VIIB:5.1:proposition}{Proposition~5.1} below in the language of cocommutative
Hopf algebras and to cite results subsequently obtained in this
direction.}
Let \(G\) be a formal \(k\)-group with affine algebra
\[
  A=\mathscr A(G).
\]
Let \(A^+\) be the augmentation ideal of \(A\). For \(a\in A\),
write
\[
  \overline a=a-\varepsilon_A(a)1_A
\]
for its projection onto \(A^+\). Let \(F\) be a formal subgroup of
\(G\), defined by a closed ideal \(J\) of \(A\), and let
\[
  \pi:A\longrightarrow A/J
\]
be the projection. Write
\[
  \overline\pi:A\longrightarrow A^+/J
\]
for the composite of the projections
\[
  A\longrightarrow A^+\longrightarrow A^+/J.
\]
For every \(a\in A\), the projection of \(\Delta_A(a)\) onto
\(A\widehat\otimes A^+\) is
\[
  \Delta_A(a)-a\widehat\otimes1.
\]

By \SGARefLink{sga3:VIIB:2.4:theorem}{Theorem~2.4}, the quotient formal \(k\)-variety \(X=G/F\)
exists. Its affine algebra is
\[
\begin{aligned}
  B
  &=
  \operatorname{Eq}\!\left(
    A
    \mathrel{\mathop{\rightrightarrows}^{\tau_1}_{%
      (\operatorname{id}_A\widehat\otimes\pi)\Delta_A}}
    A\widehat\otimes_k(A/J)
  \right)
  \\
  &=
  \left\{
    a\in A\ \middle|\
    (\operatorname{id}_A\widehat\otimes\pi)\Delta_A(a)
    =
    a\widehat\otimes\pi(1)
  \right\}
  \\
  &=
  \ker\!\left(
    (\operatorname{id}_A\widehat\otimes\overline\pi)\Delta_A
  \right),
\end{aligned}
\]
where
\[
  \tau_1(a)=a\widehat\otimes\pi(1).
\]

Equivalently, \(B\) is the subalgebra of \(A\) consisting of the
elements \(\varphi\) such that
\[
  \varphi(gh)=\varphi(g)
\]
for every \(C\in\operatorname{Alf}/k\),
\(g\in G(C)\), and \(h\in F(C)\). For
\(g,g'\in G(C)\) and \(h\in F(C)\), one has
\[
  \varphi(gg'h)=\varphi(gg').
\]
Thus \(\Delta_A(\varphi)\) belongs to the kernel of
\[
  \operatorname{id}_A\widehat\otimes_k
  \bigl((\operatorname{id}_A\widehat\otimes\overline\pi)
        \Delta_A\bigr),
\]
which equals \(A\widehat\otimes_kB\), since \(A\) is
topologically flat over \(k\). Consequently,
\[
  \Delta_A(B)\subset A\widehat\otimes_kB;
\]
that is, the closed subalgebra \(B\) is also a left coideal.

Conversely, \(B\) determines \(F\), because \SGARefLink{sga3:VIIB:2.4.1:corollary}{Corollary~2.4.1} gives
\[
  J=AB^+,
  \qquad
  B^+=B\cap A^+;
\]
in other words, \(J\) is the closed ideal generated by \(B^+\).
We therefore obtain an injective map
\[
  \mathcal Q:\mathcal F\longrightarrow\mathcal B
\]
from the set \(\mathcal F\) of formal subgroups of \(G\) to the set
\(\mathcal B\) of closed subalgebras of \(A\) that are left
coideals. \SGARefLink{sga3:VIIB:5.1:proposition}{Proposition~5.1} below shows that \(\mathcal Q\) is
bijective when \(G\) is infinitesimal. In fact, it is bijective for
every formal \(k\)-group \(G\).

Recall from \SGARefLink{sga3:VIIB:2.2.1}{2.2.1} that the functor
\[
  G\longmapsto\mathbf H(G)
\]
is an equivalence from formal \(k\)-groups to cocommutative
\(k\)-Hopf algebras. If \(F\) is defined by the closed ideal
\(J\subset A\), then the Hopf subalgebra
\[
  \mathbf H(F)\subset H=\mathbf H(G)
\]
is the orthogonal complement of \(J\) for the duality
\[
  A=H^*,
  \qquad
  H=A^\dagger
\]
\textup{(cf.\ \SGARefLink{sga3:VIIB:0.2.2}{0.2.2})}. Conversely, if \(B\) is a closed subalgebra
of \(A\) that is also a left coideal, then
\[
  I=B^\perp
\]
is a coideal and a left ideal of \(H\); explicitly,
\[
  \Delta_H(I)\subset I\otimes H+H\otimes I,
  \qquad
  \varepsilon_H(I)=0.
\]

Let \(\mathcal H\) be the set of Hopf subalgebras of \(H\), and
let \(\mathcal I\) be the set of left ideals of \(H\) that are
also coideals. For \(I\in\mathcal I\), denote by
\[
  \pi_I:H\longrightarrow H/I
\]
the projection, and by
\[
  \overline\pi_I:H\longrightarrow H^+/I
\]
the composite
\[
  H\longrightarrow H^+\longrightarrow H^+/I,
\]
where \(H^+\) is the augmentation ideal of \(H\).

Let \(K\) be a Hopf subalgebra of \(H\), and put
\[
  K^+=K\cap H^+.
\]
If \(F\) is the formal subgroup corresponding to \(K\), then
\[
  J=K^\perp
\]
and \(A^+/J\) identifies with the dual of \(K^+\). Since
\(B=\mathcal Q(F)\) is the kernel of
\[
  A
  \xrightarrow{\ \Delta_A\ }
  A\widehat\otimes_kA
  \xrightarrow{\ \operatorname{id}\widehat\otimes\overline\pi\ }
  A\widehat\otimes_k(A^+/J),
\]
duality identifies \(\mathcal Q\) with the map
\[
  \Phi:\mathcal H\longrightarrow\mathcal I,
  \qquad
  K\longmapsto HK^+.
\]
Indeed, \(HK^+\) is the image of
\[
  H\otimes_kK^+
  \xrightarrow{\ \operatorname{id}\otimes\operatorname{can}\ }
  H\otimes_kH
  \xrightarrow{\ m_H\ }
  H,
\]
and is both a left ideal and a coideal.

Similarly, the map that assigns \(AB^+\) to
\(B\in\mathcal B\) corresponds by duality to
\[
  \Psi:\mathcal I\longrightarrow\mathcal H,
\]
where
\[
  \Psi(I)
  =
  \left\{
    x\in H\ \middle|\
    (\operatorname{id}_H\otimes\pi_I)\Delta_H(x)
    =
    x\otimes\pi_I(1)
  \right\}.
\]

\medskip
\noindent\textbf{Theorem 5.0.1.}
\phantomsection
\label{sga3:VIIB:5.0.1}
{\itshape Let \(k\) be a field and \(H\) a cocommutative Hopf
algebra. The maps \(\Phi\) and \(\Psi\) above are mutually inverse
bijections between the Hopf subalgebras of \(H\) and the left
ideals of \(H\) that are coideals.}

This theorem was first proved by K.~Newman
\textup{([Ne75], \SGARefLink{sga3:VIIB:4.1}{Theorem~4.1}, where the word
``cocommutative'' was omitted)}. Newman's proof uses the
Cartier--Gabriel--Kostant theorem \textup{(cf.\ \SGARefLink{sga3:VIIB:2.9}{2.9})} to reduce to
the case in which \(H\) is connected, and then uses the existence
of a ``Sweedler basis'' in that case \textup{([Sw67],
Theorem~3)}, a result dual to the Dieudonn\'e--Cartier
\SGARefLink{sga3:VIIB:5.2.2}{Theorem~5.2.2} below. A second, shorter proof was given by
H.~J.~Schneider \textup{([Sch90], Theorem~4.15)}. A.~Masuoka
subsequently generalized the result to the case in which only the
coradical \(H_0\) of \(H\), the sum of its simple subcoalgebras, is
assumed commutative \textup{([Ma91], Theorem~1.3\,(3))}.

Finally, for a commutative \(k\)-Hopf algebra corresponding to an
affine \(k\)-group scheme \(G\), one cannot expect an analogue of
\SGARefLink{sga3:VIIB:5.0.1}{Theorem~5.0.1} without additional hypotheses: for a
\(k\)-subgroup \(F\subset G\), the quotient \(G/F\) need not be
affine. Nevertheless, M.~Takeuchi established analogous
bijections. In \textup{[Tak72], \SGARefLink{sga3:VIIB:4.3}{Theorem~4.3}}, invariant
\(k\)-subgroups \(F\) correspond to subalgebras
\(B\subset\mathcal O(G)\) satisfying
\[
  \Delta(B)\subset A\otimes B
\]
and stable under the antipode. In \textup{[Tak79], Theorem~3},
the \(k\)-subgroups \(F\) for which \(G/F\) is affine correspond to
such subalgebras \(B\) for which \(B\to A\) is faithfully flat.
